This section presents the empirical results of benchmarking ECAPA-TDNN and RawNet3 across TidyVoice and ViMD, detailing training dynamics, verification accuracy, and computational efficiency.

\subsection{Training Dynamics and Convergence}

Table~\ref{tab:training-dynamics} records the training progression across the four model--dataset pairs. Each run operated under an early stopping patience of 1 epoch on validation EER (with minDCF as tie-breaker).

\begin{table}[H]
\centering
\resizebox{\columnwidth}{!}{%
\begin{tabular}{lcrrrl}
\toprule
\textbf{Model / Dataset} & \textbf{Epoch} & \textbf{Train Loss} & \textbf{Val EER (\%)} & \textbf{Val minDCF} & \textbf{Outcome} \\
\midrule
\multirow{2}{*}{ECAPA-TDNN / TidyVoice} 
  & 0 & 16.4067 & \textbf{3.943} & \textbf{0.3685} & \textbf{Best Checkpoint} \\
  & 1 & 15.7529 & 4.075 & 0.4068 & Early Stopped \\
\midrule
\multirow{2}{*}{ECAPA-TDNN / ViMD} 
  & 0 & 17.0806 & \textbf{3.936} & \textbf{0.4503} & \textbf{Best Checkpoint} \\
  & 1 & 15.3398 & 6.634 & 0.7661 & Early Stopped \\
\midrule
\multirow{2}{*}{RawNet3 / TidyVoice} 
  & 0 & 15.7133 & \textbf{5.860} & \textbf{0.5251} & \textbf{Best Checkpoint} \\
  & 1 & 14.7825 & 9.693 & 0.7767 & Early Stopped \\
\midrule
\multirow{3}{*}{RawNet3 / ViMD} 
  & 0 & 15.9177 & 29.664 & 0.9164 & Initialized \\
  & 1 & 14.8658 & 7.532  & 0.6805 & Improved \\
  & 2 & 13.3438 & \textbf{2.692} & \textbf{0.2446} & \textbf{Best Checkpoint} \\
\bottomrule
\end{tabular}%
}
\caption{Training progression and checkpoint selection across four experimental runs.}
\label{tab:training-dynamics}
\end{table}

A notable observation is that decreasing training loss did not consistently translate into lower validation error. For ECAPA-TDNN on both datasets and RawNet3 on TidyVoice, the models achieved their best validation metrics at Epoch~0; subsequent epochs caused validation degradation, triggering early stopping. In contrast, RawNet3 on ViMD exhibited continuous and substantial convergence, reducing validation EER from 29.66\% at Epoch~0 to 2.69\% at Epoch~2.

\subsection{Speaker Verification Benchmark}

Table~\ref{tab:test-metrics} reports the primary evaluation metrics on the unseen Test splits using the best checkpoints selected on Validation.

\begin{table}[H]
\centering
\resizebox{\columnwidth}{!}{%
\begin{tabular}{llrrrrrr}
\toprule
\textbf{Model} & \textbf{Dataset} & \textbf{EER (\%)} $\downarrow$ & \textbf{minDCF} $\downarrow$ & \textbf{Accuracy (\%)} & \textbf{FAR (\%)} & \textbf{FRR (\%)} & \textbf{TAR (\%)} $\uparrow$ \\
\midrule
ECAPA-TDNN & TidyVoice & \textbf{3.912} & \textbf{0.3586} & 97.865 & 0.080 & 28.159 & \textbf{71.841} \\
RawNet3    & TidyVoice & 5.875          & 0.5624          & 96.509 & 0.124 & 46.126 & 53.874 \\
\midrule
ECAPA-TDNN & ViMD      & 3.804          & 0.4516          & 99.527 & 0.116 & 34.741 & 65.259 \\
RawNet3    & ViMD      & \textbf{3.071} & \textbf{0.2262} & \textbf{99.765} & \textbf{0.075} & \textbf{15.547} & \textbf{84.453} \\
\bottomrule
\end{tabular}%
}
\caption{Final test performance on TidyVoice and ViMD. Thresholds are calibrated at $\text{FAR} = 0.1\%$ on the Validation split and evaluated once on the Test split.}
\label{tab:test-metrics}
\end{table}

\textbf{Analysis of Operating Performance.} High overall accuracy values ($>96\%$) reflect the protocol's 100{,}000 impostor trials relative to genuine trials. Therefore, EER, minDCF, and True Acceptance Rate ($\text{TAR}$) at low False Acceptance Rates serve as the discriminative metrics:
\begin{itemize}
    \item \textbf{On TidyVoice (Multilingual):} ECAPA-TDNN achieves superior performance across all metrics, reducing EER by 33.4\% relative (3.912\% vs.\ 5.875\%) and minDCF by 36.2\% relative (0.3586 vs.\ 0.5624), while yielding a +17.97 percentage points higher TAR at the frozen security threshold.
    \item \textbf{On ViMD (Vietnamese Multi-Dialect):} RawNet3 outperforms ECAPA-TDNN, achieving a 19.3\% relative reduction in EER (3.071\% vs.\ 3.804\%) and a 49.9\% relative reduction in minDCF (0.2262 vs.\ 0.4516). At the operational threshold, RawNet3 boosts TAR to 84.453\% (compared to 65.259\% for ECAPA-TDNN) while maintaining a strict 0.075\% FAR.
\end{itemize}

Table~\ref{tab:tar-at-far} provides a fine-grained evaluation of authentication security, detailing the True Acceptance Rate at varying target False Acceptance Rates ($\text{TAR@FAR}$).

\begin{table}[H]
\centering
\resizebox{\columnwidth}{!}{%
\begin{tabular}{llrrrr}
\toprule
\textbf{Model} & \textbf{Dataset} & \textbf{TAR@FAR 5\%} & \textbf{TAR@FAR 1\%} & \textbf{TAR@FAR 0.1\%} & \textbf{TAR@FAR 0.01\%} \\
\midrule
ECAPA-TDNN & TidyVoice & \textbf{96.771} & \textbf{90.377} & \textbf{73.563} & \textbf{48.974} \\
RawNet3    & TidyVoice & 93.391          & 82.907          & 48.088          & 0.886 \\
\midrule
ECAPA-TDNN & ViMD      & \textbf{97.505} & 88.484          & 63.820          & 37.812 \\
RawNet3    & ViMD      & 97.217          & \textbf{95.106} & \textbf{86.756} & \textbf{54.798} \\
\bottomrule
\end{tabular}%
}
\caption{True Acceptance Rate (\%) at specific False Acceptance Rate thresholds ($\text{TAR@FAR}$).}
\label{tab:tar-at-far}
\end{table}

\subsection{Computational Efficiency and Model Selection}

To evaluate suitability for real-time deployment in our virtual assistant pipeline, standalone model inference latency was benchmarked on an NVIDIA Tesla T4 GPU using CUDA events across 50 iterations with batch size 1 after 10 warm-up runs (Table~\ref{tab:latency-benchmarks}).

\begin{table}[H]
\centering
\resizebox{0.9\columnwidth}{!}{%
\begin{tabular}{lrrr}
\toprule
\textbf{Model} & \textbf{Parameters} & \textbf{Median Latency (ms)} & \textbf{p95 Latency (ms)} \\
\midrule
ECAPA-TDNN & 20.77~M & 12.023 & 12.912 \\
RawNet3    & \textbf{16.28~M} & \textbf{8.840}  & \textbf{12.884} \\
\bottomrule
\end{tabular}%
}
\caption{Model complexity and isolated GPU inference latency (batch size 1 on NVIDIA Tesla T4).}
\label{tab:latency-benchmarks}
\end{table}

\textbf{Deployment Selection Rationale.} Model selection is based on Validation evidence and deployment constraints; Test results are used only for final confirmation. \textbf{RawNet3 fine-tuned on ViMD (Epoch~2)} is selected as the verification engine for the Who-Speak-AI virtual assistant:

\begin{enumerate}
    \item \textbf{Validation Evidence and Target-Domain Alignment:} The intended application targets Vietnamese speech, making ViMD the priority domain. On ViMD Validation, RawNet3 achieved 2.692\% EER and 0.2446 minDCF, outperforming ECAPA-TDNN at 3.936\% EER and 0.4503 minDCF.

    \item \textbf{Final Test Confirmation:} With the Validation-selected checkpoint and frozen Validation threshold, ViMD Test confirms the same ranking: RawNet3 achieved 3.071\% EER, 0.2262 minDCF, 0.075\% FAR, and 84.453\% TAR.

    \item \textbf{Efficiency Profile:} RawNet3 contains 21.6\% fewer parameters (16.28M vs.\ 20.77M) and exhibits 26.5\% lower median GPU latency (8.84~ms vs.\ 12.02~ms), preserving critical compute headroom for concurrent ASR, LLM reasoning, and TTS synthesis.
\end{enumerate}

\subsection{Limitations}

The experiment face multiple limitations: single-seed training (42), three epochs with early stopping, deterministic utterance per speaker per epoch; No hyperparameter grid search, margin ablation, scheduler ablation, augmentation study, multi-seed variance estimate, or CPU benchmark was performed.

The reported latency is model-only latency on a Tesla T4 and excludes audio decoding, resampling, application logic, networking, ASR, and TTS. The security threshold was calibrated on ViMD rather than live application microphones and therefore serves only as an evidence-based initial value.

WavLM+MHFA was implemented and attempted but excluded from the final comparison because no complete accepted run passed the evidence validator: one attempt produced non-finite FP16 training embeddings, and a later attempt produced non-finite evidence during ViMD Final Test. No WavLM checkpoint or metric is reported as an accepted result.
